\documentclass[a4paper,10pt]{amsart}

\pdfinfoomitdate=1
\pdftrailerid{}
\pdfsuppressptexinfo=15

\usepackage[T1]{fontenc}
\usepackage{lmodern}
\usepackage[margin=27mm]{geometry}
\usepackage{amsmath,amssymb,mathtools,booktabs,array,microtype,xcolor}
\usepackage[numbers,sort&compress]{natbib}
\usepackage[colorlinks=true,linkcolor=blue!55!black,citecolor=blue!55!black,urlcolor=blue!55!black]{hyperref}
\hypersetup{pdftitle={},pdfauthor={},pdfsubject={},pdfkeywords={}}

\newtheorem{theorem}{Theorem}
\newtheorem{proposition}[theorem]{Proposition}
\newtheorem{lemma}[theorem]{Lemma}
\newtheorem{corollary}[theorem]{Corollary}
\theoremstyle{remark}
\newtheorem{remark}[theorem]{Remark}

\newcommand{\Prob}{\mathbb P}
\newcommand{\E}{\mathbb E}
\newcommand{\one}{\mathbf 1}
\newcommand{\cS}{\mathcal S}

\title[First passage on unequal finite spiders]{Leaf-Marked First Passage on Unequal Finite Spiders:\\
A Rational Transform, Sharp Extremizers, and a Coarse-Data Inverse}
\author{Anonymous}
\date{}

\begin{document}

\begin{abstract}
Join $r\ge2$ finite paths of positive integer lengths at one centre, start
simple random walk at the centre, and absorb it at the first leaf.  We derive
an explicit continuant factorization of the probability generating functions
jointly marked by the absorbing leaf for arbitrary unequal arms.  A
centre-excursion renewal specializes the generic time/place law to a common
rational denominator for all leaf marks.  Besides parity and the first
possible atom, the specialization yields the compact formula
\[
 \operatorname{Var}(T)=\frac{C-2L}{3H}+\frac{L^2}{3H^2},
 \qquad
 H=\sum_i\ell_i^{-1},\quad L=\sum_i\ell_i,\quad C=\sum_i\ell_i^3.
\]
For fixed arm count and total length, strict integer transfers identify the
exact mean-minimizing profile---one long arm and all other arms of length
one---and the exact mean-maximizing balanced profiles.  The labelled endpoint
law determines precisely the primitive arm-length ratios and is blind to a
common dilation; adding the mean recovers that scale uniquely.  General-tree
endpoint and mean formulas, an unequal-arm endpoint exercise, equal-arm star
laws, spider spectral methods, generic finite-chain time/place laws, general
tree hitting-time generating functions, and inverse first-passage frameworks
are treated as prior background.  An exact-arithmetic audit performs
$1{,}446{,}432$ exact recurrence, coefficient, moment, extremal, and inverse
checks.
\end{abstract}

\maketitle

\section{Process, ownership boundary, and theorem}\label{sec:setup}

For positive integers $\ell_1,\ldots,\ell_r$, let
$\cS(\ell_1,\ldots,\ell_r)$ be the tree formed by identifying one endpoint
of $r$ paths of edge lengths $\ell_i$.  Call the identified vertex the
\emph{centre}, label the other endpoints $1,\ldots,r$, and make those leaves
absorbing.  Simple random walk starts at the centre.  Write $T$ for its first
leaf-hitting time and $I$ for the label of the leaf it hits.

Unequal arms obstruct the radial lumping available when all lengths agree:
the time and the absorbing label remain coupled.  The paper resolves that
coupling through a leaf-marked excursion transform.  The contribution is a
four-part conjunction: the explicit unequal-spider continuant factorization,
its compact scalar variance specialization, sharp fixed-total extremizers,
and a geometry inverse from deliberately coarse data.  Generic existence,
rationality, resolvent, time/place, and second-moment formulas receive zero
credit.  Table~\ref{tab:subtraction} states the corresponding subtraction.

\begin{table}[ht]
\centering
\small
\caption{Primary-source subtraction.  The middle column receives zero
contribution credit; the right column is the narrower residual studied here.}
\label{tab:subtraction}
\begin{tabular}{@{}>{\raggedright\arraybackslash}p{.23\textwidth}
                    >{\raggedright\arraybackslash}p{.335\textwidth}
                    >{\raggedright\arraybackslash}p{.335\textwidth}@{}}
\toprule
Source & Owned input & Boundary retained here \\
\midrule
Pearce~\cite{Pearce1980} & finite trees with absorbing leaves; endpoint
probabilities and expected walk length & neither endpoint mass nor mean is
claimed \\
Pal--Mesikepp, Problem 2.4~\cite{PalMesikepp2025} & designated-leaf
probability for a star with unequal integer arms & endpoint law is used only
for extrema and inversion \\
Castella--Sericola~\cite{CastellaSericola2026} & hitting distributions and
moments for equal-length finite stars & the transform axis requires arbitrary
unequal arms and an explicit continuant factorization \\
Sericola~\cite{Sericola2024} & generic finite-chain joint hitting-time/place
law and first/second moment matrices & only the closed unequal-spider
continuant product and compact scalar specialization \\
Chen~\cite{Chen2007} & algorithms for hitting-time generating functions on
general trees & a leaf-marked, unequal-spider closed factorization rather
than generic existence or an algorithm \\
de la Iglesia--Juarez~\cite{DelaIglesiaJuarez2023} & spectral and stochastic
factorization methods for half-line spider chains & no spectral-framework or
generic-resolvent claim \\
de la Pe{\~n}a--Gzyl--McDonald~\cite{DelaPenaGzylMcDonald2008} & unknown
transitions on a known augmented tree recovered from rich boundary time/place
data & fixed simple-walk kernel and unknown integer arms, using only endpoint
probabilities plus one mean \\
\bottomrule
\end{tabular}
\end{table}

Sericola's generic identity
$\Prob_i(T_A=n,X_{T_A}=j)=(P_{A^c}^{n-1}P_{A^c,A})_{ij}$ already owns the
joint time/place law and its matrix moments~\cite{Sericola2024}.  Chen's
general-tree work already owns algorithmic hitting-time generating
functions~\cite{Chen2007}.  Neither inspected source states the product
formula \eqref{eq:marked} for unequal finite spiders or its continuant
denominator; that explicit specialization is the only residual transform
claim.

The term ``reflecting--absorbing factorization'' in
\cite{DelaIglesiaJuarez2023} denotes a stochastic block factorization and a
Darboux transformation, not absorption at finite leaves.  Conversely, the
inverse theorem in \cite{DelaPenaGzylMcDonald2008} observes complete joint
time/place laws on two boundary layers and recovers transition probabilities;
our topology class and transition rule are fixed, while the integer arm
lengths are unknown.  These distinctions prevent generic spider, resolvent,
and tomography language from carrying contribution credit.

Define continuant polynomials by
\begin{equation}\label{eq:P-rec}
 P_0(z)=0,\qquad P_1(z)=1,\qquad P_2(z)=2,
 \qquad P_m(z)=2P_{m-1}(z)-z^2P_{m-2}(z).
\end{equation}
Put
\begin{equation}\label{eq:D-def}
 P(z)=\prod_{j=1}^rP_{\ell_j}(z),\qquad
 D(z)=rP(z)-z^2\sum_{i=1}^rP_{\ell_i-1}(z)
                 \prod_{j\ne i}P_{\ell_j}(z).
\end{equation}
Also write
\begin{equation}\label{eq:HLC}
 H=\sum_{i=1}^r\frac1{\ell_i},\qquad
 L=\sum_{i=1}^r\ell_i,\qquad C=\sum_{i=1}^r\ell_i^3.
\end{equation}

\begin{theorem}[marked law, moments, extrema, and inverse]\label{thm:main}
For the walk on $\cS(\ell_1,\ldots,\ell_r)$, the following statements hold.
\begin{enumerate}
\item For every labelled leaf $i$,
\begin{equation}\label{eq:marked}
 F_i(z):=\E\!\left[z^T\one_{\{I=i\}}\right]
 =\frac{z^{\ell_i}\prod_{j\ne i}P_{\ell_j}(z)}{D(z)}.
\end{equation}
Its support lies in $\ell_i+2\mathbb Z_{\ge0}$, and
\begin{equation}\label{eq:first-atom}
 \Prob(T=\ell_i,I=i)=\frac1{r2^{\ell_i-1}}.
\end{equation}
\item The endpoint and mean identities
\begin{equation}\label{eq:background-moments}
 \Prob(I=i)=\frac{\ell_i^{-1}}H,\qquad \E T=\frac LH
\end{equation}
are background inputs from the general-tree and unequal-arm literature in
Table~\ref{tab:subtraction}.  The residual second-moment conclusion is
\begin{equation}\label{eq:variance}
 \operatorname{Var}(T)=\frac{C-2L}{3H}+\frac{L^2}{3H^2}.
\end{equation}
\item Fix $r$ and $L\ge r$, and write $L=qr+s$ with $0\le s<r$.  Then
\begin{equation}\label{eq:extremal}
 \frac{L}{r-1+1/(L-r+1)}\le \E T
 \le \frac{L}{(r-s)/q+s/(q+1)}.
\end{equation}
Equality on the left holds exactly for permutations of
$(L-r+1,1,\ldots,1)$; equality on the right holds exactly when the arms are
balanced, with $r-s$ lengths $q$ and $s$ lengths $q+1$.
\item The labelled endpoint vector determines the primitive positive integer
arm vector $d=(d_i)$, while $\ell=cd$ is invisible for every common positive
integer dilation $c$.  Adding $\E T$ determines the unique scale by
\begin{equation}\label{eq:scale}
 c^2=\E T\,\frac{\sum_i d_i^{-1}}{\sum_i d_i}.
\end{equation}
\end{enumerate}
\end{theorem}

The theorem assumes from the outset that the unknown object is a labelled
finite spider and that its transition kernel is simple random walk.  Part~(4)
does not recover an arbitrary tree or unknown transition probabilities.

\section{Killed paths and centre renewal}\label{sec:transform}

Let $U_m$ denote the Chebyshev polynomial of the second kind, with
$U_{-1}=0$.  For $m\ge1$, recurrence \eqref{eq:P-rec} is equivalently the
polynomial identity
\begin{equation}\label{eq:Cheb}
 P_m(z)=z^{m-1}U_{m-1}(1/z).
\end{equation}
The Chebyshev notation is a convenient solution of a second-order recurrence;
neither generic continuants nor gambler's ruin are contribution claims.

\begin{lemma}[one killed arm]\label{lem:arm}
On $\{0,1,\ldots,m\}$, start fair nearest-neighbour walk at $1$ and stop on
first hitting $\{0,m\}$.  If $\tau$ is this path time, then
\begin{align}
 \E_1[z^\tau\one_{\{X_\tau=m\}}]
   &=\frac1{U_{m-1}(1/z)}=\frac{z^{m-1}}{P_m(z)},\label{eq:success-arm}\\
 \E_1[z^\tau\one_{\{X_\tau=0\}}]
   &=\frac{U_{m-2}(1/z)}{U_{m-1}(1/z)}
     =\frac{zP_{m-1}(z)}{P_m(z)}.\label{eq:return-arm}
\end{align}
\end{lemma}

\begin{proof}
For $1\le k\le m-1$, either boundary-marked transform $u_k$ satisfies
\[
 u_k=\frac z2(u_{k-1}+u_{k+1}).
\]
For absorption at $m$, the boundary values are $(u_0,u_m)=(0,1)$, and the
solution is
\[
 u_k=\frac{U_{k-1}(1/z)}{U_{m-1}(1/z)}.
\]
For absorption at zero, the boundary values are $(1,0)$, and the solution is
\[
 u_k=\frac{U_{m-k-1}(1/z)}{U_{m-1}(1/z)}.
\]
Evaluating at $k=1$ and applying
\eqref{eq:Cheb} proves both identities.  For $m=1$, the formulas use
$P_0=0$ and express immediate success after the later-added centre step.
\end{proof}

\begin{proof}[Proof of Theorem~\ref{thm:main}(1)]
An attempt from the centre chooses arm $j$, makes one step to its first
vertex, and then either reaches leaf $j$ or returns to the centre.  By
Lemma~\ref{lem:arm}, the transform of a successful attempt at leaf $i$ is
\[
 S_i(z)=\frac1r\frac{z^{\ell_i}}{P_{\ell_i}(z)},
\]
and the transform of an unsuccessful attempt, with the arm unmarked, is
\[
 Q(z)=\frac{z^2}{r}\sum_{j=1}^r
            \frac{P_{\ell_j-1}(z)}{P_{\ell_j}(z)}.
\]
After every failure the walk is again at the centre, so the strong Markov
property gives the formal renewal series
\[
 F_i(z)=S_i(z)\sum_{k\ge0}Q(z)^k=\frac{S_i(z)}{1-Q(z)}.
\]
Here $Q(0)=0$, so the inverse is valid as a formal series.  Also
$P_m(1)=m$, whence
\[
 Q(1)=1-\frac Hr<1,
 \qquad D(1)=H\prod_{j=1}^r\ell_j>0.
\]
Thus the rational identity is regular at $z=1$ and its endpoint and moment
evaluations below are justified analytically as well as probabilistically.
Multiplying numerator and denominator by $r\prod_jP_{\ell_j}(z)$ gives
\eqref{eq:marked}.

Each $P_m$ is a polynomial in $z^2$.  Equation \eqref{eq:marked} therefore
has only powers congruent to $\ell_i$ modulo two, in agreement with
bipartiteness.  At the first possible time $\ell_i$, the walk must choose arm
$i$ and then follow its unique shortest path without backtracking.  The
probability is $(1/r)(1/2)^{\ell_i-1}$, proving \eqref{eq:first-atom}.
\end{proof}

\section{Excursion moments and the variance}\label{sec:moments}

The renewal proof also separates the second moment into one-arm quantities.
This route avoids treating differentiation of the common denominator as a
black box.

\begin{lemma}[one-attempt moments]\label{lem:attempt}
Condition on choosing an arm of length $m$.  Let $A_m$ be the duration from
departure at the centre until the attempt next reaches either the centre or
the leaf, including the departure step, and let $R_m$ be the event of return
to the centre.  Then
\begin{align}
 \Prob(R_m^c)&=\frac1m,&
 \E A_m&=m,\label{eq:attempt-first}\\
 \E A_m^2&=\frac{m(m^2+2)}3,&
 \E[A_m\one_{R_m}]&=\frac{2(m^2-1)}{3m}.
 \label{eq:attempt-second}
\end{align}
\end{lemma}

\begin{proof}
Including the centre step, the success and return transforms from
Lemma~\ref{lem:arm} are
\[
 A_m^{\rm suc}(z)=\frac{z^m}{P_m(z)},\qquad
 A_m^{\rm ret}(z)=\frac{z^2P_{m-1}(z)}{P_m(z)}.
\]
Differentiating \eqref{eq:P-rec} and inducting gives
\begin{align*}
 P_m(1)&=m,\\
 P_m'(1)&=-\frac{m(m-1)(m-2)}3,\\
 P_m''(1)&=\frac{m(m-1)(m-2)(m^2-7m+7)}{15}.
\end{align*}
Substitution into the two quotient derivatives gives
$A_m^{\rm suc}(1)=1/m$, $(A_m^{\rm ret})(1)=1-1/m$,
\[
 (A_m^{\rm suc}+A_m^{\rm ret})'(1)=m,
 \]
\[
 (A_m^{\rm suc}+A_m^{\rm ret})''(1)
 +(A_m^{\rm suc}+A_m^{\rm ret})'(1)=\frac{m(m^2+2)}3,
\]
and $(A_m^{\rm ret})'(1)=2(m^2-1)/(3m)$.  These are precisely
\eqref{eq:attempt-first}--\eqref{eq:attempt-second}.  The identities remain
valid at $m=1$.
\end{proof}

\begin{proof}[Proof of Theorem~\ref{thm:main}(2)]
Average one attempt over the uniform choice of arm.  Let $B$ indicate a
successful attempt and $A$ its duration.  Lemma~\ref{lem:attempt} gives
\begin{equation}\label{eq:renewal-data}
 p:=\E B=\frac Hr,\quad
 \mu:=\E A=\frac Lr,\quad
 \nu:=\E A^2=\frac{C+2L}{3r},\quad
 \rho:=\E[A(1-B)]=\frac{2(L-H)}{3r}.
\end{equation}
Let $T'$ be an independent copy of $T$, used after a failed attempt.  Then
$T=A+(1-B)T'$.  First moments give $p\E T=\mu$, which recovers the prior mean
in \eqref{eq:background-moments}.  Squaring before taking expectations gives
\[
 p\E T^2=\nu+2\rho\E T.
\]
Using \eqref{eq:renewal-data} and subtracting $(L/H)^2$ yields
\eqref{eq:variance}.  Finally, evaluating \eqref{eq:marked} at $z=1$ gives
the prior endpoint identity in \eqref{eq:background-moments}.
\end{proof}

\section{Sharp fixed-mass extrema}\label{sec:extrema}

The prior mean identity becomes a residual result when optimized with exact
integer equality classes.

\begin{proof}[Proof of Theorem~\ref{thm:main}(3)]
Because $L$ is fixed, minimizing $\E T=L/H$ is equivalent to maximizing
$H=\sum_i1/\ell_i$.  If $2\le a\le b$, the outward unit transfer
$(a,b)\mapsto(a-1,b+1)$ changes the reciprocal sum by
\[
 \frac1{a(a-1)}-\frac1{b(b+1)}>0.
\]
Repeated outward transfers end only when all but one arm have length one.
Strictness proves that the maximizers of $H$ are exactly the permutations of
$(L-r+1,1,\ldots,1)$.  Substitution gives the left side of
\eqref{eq:extremal}.

For the other direction, if $b\ge a+2$, the inward transfer
$(a,b)\mapsto(a+1,b-1)$ changes the old reciprocal sum minus the new one by
\[
 \frac1{a(a+1)}-\frac1{b(b-1)}>0.
\]
Thus $H$ decreases strictly until all arm lengths differ by at most one.
Writing $L=qr+s$ forces $r-s$ lengths $q$ and $s$ lengths $q+1$; substituting
their reciprocal sum gives the right side of \eqref{eq:extremal}.  These
strict transfers also prove both equality classifications.  If $L=r$, both
classes reduce to the all-one profile, so the boundary case is included.
\end{proof}

\section{What the coarse data identify}\label{sec:inverse}

\begin{proof}[Proof of Theorem~\ref{thm:main}(4)]
Let $\pi_i=\Prob(I=i)$.  The background endpoint identity gives
\begin{equation}\label{eq:ratio}
 \frac{\pi_i}{\pi_j}=\frac{\ell_j}{\ell_i}.
\end{equation}
Hence the labelled vector $\pi$ determines the rational ray containing
$(\ell_i)$ and therefore its unique primitive positive integer representative
$d=(d_i)$.  It cannot determine scale: every $cd$, $c\in\mathbb Z_{>0}$,
has the same endpoint vector.  Conversely, if two labelled spiders have the
same endpoint vector, \eqref{eq:ratio} makes their arm vectors proportional;
primitive integer normalization shows that common dilation is the entire
ambiguity.

For $\ell=cd$, the prior mean formula becomes
\[
 \E T=c^2\frac{\sum_i d_i}{\sum_i d_i^{-1}}.
\]
Solving gives \eqref{eq:scale}.  For data generated by an integer spider, the
right-hand side is the square of its unique positive integer scale.  No claim
is made for arbitrary noisy vectors, unknown topology, or unknown transition
kernels.
\end{proof}

\section{Exact audit, limitations, and conclusion}\label{sec:audit}

A standalone standard-library verifier checks the formulas in exact integer
and rational arithmetic.  For every ordered profile with $2\le r\le5$ and
$1\le\ell_i\le4$, it compares a literal vertex-state Markov recursion with
the rational series \eqref{eq:marked} coefficient by coefficient through time
$2L+12$.  Additional exact routes check quotient derivatives, the excursion
moments, every positive composition with $2\le r\le6$ and $L\le24$ for the
sharp bounds, primitive-ray recovery through arm length eight, and the
equal-arm radial collapse as an owned control.  Table~\ref{tab:audit} records
the frozen run.

\begin{table}[ht]
\centering
\small
\caption{Exact-arithmetic falsification audit.  Counts are assertions, not
proof steps or evidence of novelty.}
\label{tab:audit}
\begin{tabular}{@{}lr@{}}
\toprule
Control & Assertions \\
\midrule
Continuant recurrence and closed coefficients & 400 \\
Literal marked coefficients and mass & 524,816 \\
Endpoint/mean/variance derivatives & 10,448 \\
Additional one-attempt moments & 400 \\
Fixed-mass inequalities and equality classes & 760,104 \\
Primitive ray, dilation, and scale recovery & 149,760 \\
Equal-arm owned-background collapse & 504 \\
\midrule
Total & 1,446,432 \\
\bottomrule
\end{tabular}
\end{table}

The results have four deliberate limitations.  First, the topology class is
known and consists only of labelled spiders with positive integer arm lengths.
Second, transitions are the unweighted simple-walk transitions; weighted or
biased arms are outside the theorem.  Third, the inverse uses exact population
quantities and gives no estimator, stability bound, or noisy-data guarantee.
Fourth, the paper does not address cover times, general trees, generic
time/place laws, or equal-arm distribution theory.  Sericola's generic joint
law and moment matrices and Chen's general-tree PGF algorithm are zero-credit
inputs.  A bounded primary-source search did not locate the specific
continuant-factorization/variance/extremal/inverse conjunction, but a search
non-hit is not a novelty,
priority, or freedom-to-operate certificate.

Within those boundaries, centre renewal gives an explicit unequal-spider
continuant factorization of the generic marked law, while the prior endpoint
and mean identities
support two additional exact questions: which integer geometries extremize
the clock, and which geometry information survives coarse observation.  The
answers are respectively the strict unbalanced/balanced equality classes and
the exact common-dilation boundary.

\section*{Declarations}

\paragraph{Data availability.}
No external dataset is used.  The paper-local exact verifier and its frozen
text output accompany this anonymous internal manuscript.

\paragraph{Ethics statement.}
The work is mathematical and uses no human participants, personal data,
animals, or interventions.

\paragraph{Author contributions.}
The anonymous author or authors performed the conceptualization, derivations,
exact verification, source audit, and manuscript preparation.  Identity and
individual role allocation remain suppressed for anonymous review.

\paragraph{Conflict of interest.}
No conflict of interest is declared.

\paragraph{Funding.}
No external funding is declared for this manuscript.

\paragraph{External status.}
The manuscript remains under \texttt{HOLD\_EXTERNAL}; compilation does not
authorize posting, submission, circulation, or author contact.

\bibliographystyle{plainnat}
\bibliography{references}

\end{document}
